\section{What successive ideas achieve at concrete sizes}\label{sec:numerical}
The following tables use exactly the same values of \(n\) throughout.
They count \emph{orders}, not voters. An upper-bound entry need not
be attained by a domain. Our small constructed examples illustrate
the counts; these tables limit \emph{every} domain.

\subsection{First compare the arguments without classification data}
Use only \(f(3)=4\) as the starting value. There are three stages:
\begin{enumerate}[label=\arabic*.]
\item \textbf{Deletion alone.} The weak count \(M\leq\Sigma\)
gives \(f(n)\leq nf(n-1)\), hence \(f(n)\leq4n!/6\).
\item \textbf{The gap argument.} Discarding \(G\geq0\) from
\(G+(n+1)M\leq3\Sigma\) gives
\(f(n)\leq3n f(n-1)/(n+1)\).
\item \textbf{The full argument.} The local estimate and components
also give \(f(n)\leq32n f(n-1)/(11n+8)\). At each step use
the stronger of this and the gap-only step.
\end{enumerate}
Round down after \emph{each} recurrence step, since domain sizes
are integers. Precisely, set \(u_3=v_3=4\) and, for \(n\geq4\),
\[
 u_n=\left\lfloor\frac{3n}{n+1}u_{n-1}\right\rfloor,\qquad
 v_n=\min\!\left\{
 \left\lfloor\frac{3n}{n+1}v_{n-1}\right\rfloor,
 \left\lfloor\frac{32n}{11n+8}v_{n-1}\right\rfloor\right\}.
\]
Induction proves \(f(n)\leq u_n\) and \(f(n)\leq v_n\).
These retain the finite correction in the recurrence, so are sharper
than simply substituting \(n\) into \(4(32/11)^{n-3}\).

\begin{center}\small
\renewcommand{\arraystretch}{1.18}
\begin{tabular}{@{}rrrr@{}}\toprule
\(n\)&Deletion alone&Gap argument&Full argument\\\midrule
5&$80$&$22$&$22$\\
6&$480$&$56$&$56$\\
7&$3\,360$&$147$&$147$\\
8&$26\,880$&$392$&$392$\\
9&$241\,920$&$1\,058$&$1\,055$\\
10&$2\,419\,200$&$2\,885$&$2\,861$\\
20&$1.622\times10^{18}$&$89\,223\,302$&$77\,257\,942$\\
30&$1.769\times10^{32}$&$3.570\times10^{12}$&$2.523\times10^{12}$\\
100&$6.222\times10^{157}$&$2.743\times10^{45}$&$3.097\times10^{44}$\\
\bottomrule\end{tabular}
\end{center}
Long entries are rounded \emph{upwards} to four significant digits,
so the displayed numbers are still valid upper bounds. The complete
integers are provided in \texttt{data/upper\_bounds\_exact.csv}.

The two recurrence factors agree at \(n=8\). The gap-only factor
is smaller for \(n<8\); the \(32/11\) factor is smaller for \(n>8\).
This explains why the last two columns agree at small \(n\).
The asymptotic improvement becomes more visible as \(n\) increases.

\subsection{What the separately known small maxima improve}
The published classifications give
\[
 f(5)=20,\qquad f(6)=45,\qquad f(7)=100,\qquad f(8)=224.
\]
The first three values come from the classification through seven
alternatives~\cite{smallcd}; the eight-alternative value is
Theorem~4.1 of~\cite{eightcd}. Use these as additional upper bounds
at their respective steps, then continue the same two recurrences.
This changes the starting information, while the proof steps stay the same.

\begin{center}\small
\renewcommand{\arraystretch}{1.18}
\begin{tabular}{@{}rrr@{}}\toprule
\(n\)&Gap argument, with classifications&Full argument, with classifications\\\midrule
5&$20$&$20$\\
6&$45$&$45$\\
7&$100$&$100$\\
8&$224$&$224$\\
9&$604$&$602$\\
10&$1\,647$&$1\,632$\\
20&$50\,935\,157$&$44\,070\,458$\\
30&$2.038\times10^{12}$&$1.439\times10^{12}$\\
100&$1.566\times10^{45}$&$1.767\times10^{44}$\\
\bottomrule\end{tabular}
\end{center}
The rows through eight now contain exact classified maxima. The
later rows remain upper estimates, and are not claims of exact maxima.
In particular, knowing \(f(8)=224\) helps the finite bounds: for
\(n\geq8\), even the simple geometric consequence becomes
\[
 f(n)\leq224(32/11)^{n-8}.
\]
It improves the leading constant, rather than the exponential base.
Thus the eight-alternative classification is useful additional
information, while the proof of the base \(32/11\) does not require it.

\outside{The recurrences used here are consequences of the formal
gap and transfer results and the deletion upper estimate. These
tables are separate exact-arithmetic calculations, not new exported
Lean theorems. \texttt{scripts/build\_numerical\_examples.py} regenerates
their exact entries, checks the small constructed domains and rounds
long entries upwards. The published classifications in the second
table are external inputs; they have not been added to the assumptions
or dependency chain of the Lean proof.}
